%% =====================================================================
%%  T-23-01  Reading a Seeded Room
%%  -------------------------------------------------------------------
%%  One idea per file. Core is mandatory. Slots in canonical order.
%%  Max 700 words. No layout commands. No filler. New source material
%%  for this entry goes in sources/<BOOK>/T-23-01.tex -- never by
%%  rewriting this file (docs/RULES.md R0.1).
%% =====================================================================
%% @kind: topic
%% @id: T-23-01
%% @title: Reading a Seeded Room
%% @subtitle: Detecting consensus that has been assembled
%% @chapter: c23
%% @order: 10
%% @status: stable
%% @sources: B4
%% @updated: 2026-09-19
%% @allow-rewrite: slot migration to the seven-question format (docs/RULES.md section 2)

\topic{T-23-01}{Reading a Seeded Room}{Detecting consensus that has been assembled}


\begin{core}
A room in which everyone agrees is evidence of two possible things: that the proposition is true, or that the room was built. The two are distinguished by asking questions about the room rather than about the proposition.
\end{core}

\begin{mechanism}
Direct evaluation of the claim is useless here, because the claim is not what has been engineered. What has been engineered is the distribution of visible agreement, so the diagnostic has to examine that distribution.

Genuine consensus has a particular shape. It is uneven --- some people agree strongly, some weakly, some for different reasons. It is stable across venues and across time. It resolves into named individuals who can be contacted separately. And it tolerates dissent, because people who hold a view for reasons do not need objections suppressed.

Assembled consensus fails most of those tests at once. It is unanimous, unexplained, synchronized in time, concentrated in a single venue, anonymous or newly created, and intolerant of disagreement. Any one of these is weak evidence; three together are strong.
\end{mechanism}

\begin{application}
  \item Ask for names, then contact two of them separately and out of the venue.
  \item Check the timestamps: when did agreement become visible, and did it arrive gradually or at once?
  \item Take the same question to a room nobody built --- a different audience, a different channel, a private conversation.
  \item Ask what would have to be excluded for the room to look like this.
  \item Look at the shape of the agreement. Uneven and reasoned is real; flat and unanimous is not.
\end{application}

\begin{feedback}
  \item Agreement is unanimous and nobody can give a reason for it.
  \item Support appeared at one moment rather than accumulating.
  \item The same voices appear in every venue where the view is held.
  \item Dissent is filtered out, or answered by other audience members.
  \item Supporters are anonymous, aggregated or recently created.
  \item You cannot name one holder of the view who is not part of the presentation.
\end{feedback}

\begin{failure}
The diagnostic produces false positives on things that are merely popular. Fashion, convention and genuine majority opinion all look partly like assembled consensus, and treating every agreement as staged makes you unworkable and alone. The test is worth running where the stakes are high, where the agreement appeared suddenly, or where someone benefits from it --- not as a general posture.
\end{failure}

\begin{countermeasures}
  \item When names are refused, treat the refusal as the answer rather than pursuing the claim.
  \item When the view is defended by other audience members, notice who is being protected --- the presenter who does not answer is the one with the stake.
  \item When the timestamps are unavailable, ask why the history of the agreement is not visible.
  \item When your own dissent is treated as a problem to be corrected, agree that it is and leave; a room that cannot hold an objection was not built to hold information.
\end{countermeasures}

\sources{B4}
\seesources
\seealso{T-13-01, T-15-01, T-15-02, T-23-02}
